\begin{trapbox}
\begin{description}[style=nextline, leftmargin=2.8em, labelsep=0.5em, itemsep=0.35em]
    \item[\textbf{\textcolor{CTrap}{易错点一（忽略符号判断）}}] 在使用结论时，不先判断 $a,b$ 的符号，直接套用 $x+\frac{1}{x} \geq 2$，导致在异号情况下错误地得到 $\geq 2$ 的结论。
    \item[\textbf{\textcolor{CTrap}{正确理解（符号定方向）：}}] 结论分为同号和异号两种情况，不等号方向相反。必须先判断 $ab$ 的符号，再选择对应结论。
    \item[\textbf{\textcolor{CTrap}{错因分析（记忆片面）：}}] 只记住了常见的“正数与其倒数和不小于2”，而忽略了异号情形，或者不知道异号时结论相反。
\end{description}

\medskip
\begin{description}[style=nextline, leftmargin=2.8em, labelsep=0.5em, itemsep=0.35em]
    \item[\textbf{\textcolor{CTrap}{易错点二（取等条件误解）}}] 认为等号成立条件 $a=b$ 在所有情况下都能成立，忽略了在 $ab<0$ 时 $a=b$ 不可能成立，导致错误地认为能取到等号。
    \item[\textbf{\textcolor{CTrap}{正确理解（取等需可能）：}}] 等号成立条件 $a=b$ 仅在 $ab>0$ 时有可能成立；在 $ab<0$ 时，等号无法取得，不等式严格成立。
    \item[\textbf{\textcolor{CTrap}{错因分析（条件脱节）：}}] 没有将等号成立条件与前提条件结合起来分析，孤立地记忆结论。
\end{description}

\medskip
\begin{description}[style=nextline, leftmargin=2.8em, labelsep=0.5em, itemsep=0.35em]
    \item[\textbf{\textcolor{CTrap}{易错点三（滥用最值结论）}}] 误认为 $x+\frac{1}{x}$ 在 $x<0$ 时有最小值，或者误认为 $a/b+b/a$ 在 $ab<0$ 时有最小值，从而错误求值。
    \item[\textbf{\textcolor{CTrap}{正确理解（最值有别）：}}] 当 $ab>0$ 时，$a/b+b/a$ 有最小值 $2$；当 $ab<0$ 时，$a/b+b/a$ 有最大值 $-2$，且最大值不一定能取到。
    \item[\textbf{\textcolor{CTrap}{错因分析（最值混淆）：}}] 将同号情形下的最小值结论机械地推广到异号情形，没有注意符号改变时不等号方向也改变。
\end{description}
\end{trapbox}
